\section{Data: creación de datos y construcción de capacidades}

\subsection{Nuevo Domain = nueva Capability}
El instructor insiste en que K2.5 no consiste en poner etiquetas sobre un modelo antiguo, sino en que ``un nuevo domain significa una nueva rebanada de capacidad'': imagen y texto, API, GUI, herramientas de escritorio e incluso la consola de OpenClaw se diseñaron, cada uno, como un benchmark propio. La recolección de datos, las normas de etiquetado y los success criteria de cada domain determinan qué tipo de conducta terminará aprendiendo el agente. La primera pregunta que debemos hacernos al leer el reporte técnico es: ¿quién define los datos de este domain, cómo se sintetizan y qué tipo de salida producen?

La estrategia de mezcla temprana incorporaba imágenes en menos del 10\% (1:9); a la mitad del entrenamiento, en el paso 250,000, pasó a 2:8, y al final llegó a 1:1. Pero la frase clave del instructor es: ``basta con introducir la visión desde el Step0 para que las capacidades de texto y de programación suban al mismo tiempo'', lo que muestra que la visión no es un añadido, sino una dimensión central que moldea la comprensión del lenguaje.

\begin{importantbox}{La palanca de mezclar visión desde el inicio}
En las primeras etapas del entrenamiento solo se destinaba cerca del 10\% a datos de imagen y texto (1:9); en la etapa intermedia se llegó al 20\% (2:8), y solo hacia el final se alcanzó 1:1. Los datos del instructor muestran que mezclar visión desde el Step0 mejora de manera notable las capacidades de \textbf{texto} y \textbf{programación}, lo que indica que la visión en K2.5 no es un elemento adicional, sino que inyecta continuamente una nueva señal en la representación del lenguaje.
\end{importantbox}

\subsection{Métodos de síntesis de datos}
El sistema de datos de K2.5 es altamente reproducible: se descompone la capacidad objetivo en benchmarks (por ejemplo, instrucciones en pantalla, respuestas entre dominios, razonamiento sobre textos largos), se diseña un pipeline de generación de datos para cada benchmark y luego se entrena un baseline con la combinación de SFT + RL para verificar la calidad de los datos. El instructor insiste en que debemos dejar por escrito el ``criterio de terminación'' de cada benchmark, de modo que se pueda saber en qué momento el modelo terminó la task.

El flujo habitual sigue siendo 1) definir el benchmark, 2) generar muestras con reglas o plantillas, 3) validarlas con el baseline antes de pasar a RL. Este pipeline se menciona repetidamente en la sección de la clase entre 22:00 y 32:00, sobre todo cuando, al sintetizar muestras de video-to-code o de GUI, se escriben en el prompt tanto el trace de la API como el estado de la interfaz de usuario a la vez.

\begin{enumerate}
    \item \textbf{Definir el benchmark}: especificar varios tipos de tareas (como video-to-code, operación de GUI, comprensión de textos largos) y precisar los success criteria y los indicadores visualizables.
    \item \textbf{Diseñar el pipeline algorítmico de generación}: usar recuperación de imágenes, un tool simulator y una gramática impulsada por prompts para generar muestras sintéticas, registrando los metadatos de cada paso.
    \item \textbf{Validación con baseline}: correr el modelo SFT sobre el benchmark y confirmar el puntaje de simulación del lado del reward antes de permitir el paso a RL.
\end{enumerate}

\begin{table}[H]
\footnotesize
\centering
\begin{tabular}{p{3cm}p{4cm}p{7cm}}
\toprule
Etapa de entrenamiento & Proporción visión / texto & Efecto observado \\
\midrule
Inicial (Step0) & 1:9 & Aunque las muestras visuales son pocas, basta con mezclarlas desde el primer paso para que el desempeño en texto y en código mejore con rapidez, produciendo en poco tiempo un \textbf{refuerzo del conocimiento} en varios benchmarks. \\
Intermedia (Step 250,000) & 2:8 & Al aumentar la proporción de visión, las capas intermedias aprenden a insertar llamadas a herramientas, operaciones de GUI y percepción de escena dentro de la secuencia, lo que alivia los conflictos de gradiente entre modalidades. \\
Final (etapa final) & 1:1 & Visión y texto reciben el mismo peso, de modo que los critical steps de la evaluación final examinan a la vez el razonamiento visual y la ejecución del lenguaje, mejorando el agentic performance. \\
\bottomrule
\end{tabular}
\caption{Proporción de mezcla de imagen y texto controlada por K2.5 en cada etapa y su efecto en el entrenamiento, con datos tomados de la curva de la Figura 9 del apéndice del instructor.}
\end{table}

\subsection{Domain-specific synthesis loops}
Cada benchmark tiene sus propios detalles de construcción: video-to-code requiere cortar en varios cuadros y alinear con texto la UI, la landing page, los endpoints de la API y los bocetos de diseño; la comprensión de textos largos requiere un OCR consciente del layout; las muestras de GUI/CLI deben conservar el tool trace, simular clics y registrar las respuestas de la API. El instructor insiste en que estos pipelines de datos determinan de manera directa la estructura del prompt del agente y la longitud de la generación, por lo que se convierten en el eslabón clave de ``domain-to-capability''. En la clase, entre 15:40 y 18:10, vemos que el instructor usa una tabla de dos columnas, ``Tool trace + estado de la UI'', para enumerar los campos que deben registrarse en cada sample.

\begin{knowledgebox}{Tres estrategias de domain-specific synthesis}
\begin{itemize}
    \item Descomponer el benchmark de cada domain en tres segmentos, intent, tools y output, para precisar cuándo se necesita GUI, cuándo API y cuándo basta con texto.
    \item Usar un tool simulator y un UI mock generator para producir muestras con trace, garantizando un token budget consistente durante el entrenamiento.
    \item Usar un uniform metadata schema para registrar timestamp, frame id y tool log, de modo que la later evaluation pueda reproducir cada critical step.
\end{itemize}
\end{knowledgebox}

\subsection{Gobernanza de calidad e interpretabilidad}
El instructor señala que el pipeline de datos no debe preocuparse solo por la cantidad, sino mantener la quality dentro del límite de los ``critical steps''. Cada lote de muestras se corre una vez sobre el baseline de SFT; si el puntaje de reward queda por debajo del umbral, se regresa al módulo de datos para recalibrar la prompt grammar. Además, las visual sample llevan una frame-level label adjunta, lo que facilita, en la later debugging, rastrear cada Token visual que produjo un error.

\begin{importantbox}{El Guardrail automatizado de la calidad de los datos}
K2.5 usa dos niveles de revisión de calidad: el primero es el reward gate del baseline de SFT, y el segundo es una revisión human-in-the-loop de las operaciones de GUI y del texto en pantalla. Solo las sample que pasan ambos niveles a la vez pueden entrar a RL, lo que evita el ruido explosivo del ``reward hacking''.
\end{importantbox}

\subsection{Resumen de la sección}
La historia de datos de K2.5 es un mapeo de ``domain-to-capability'': encadena la capacidad de datos a través de tres categorías de benchmark, video-to-code, razonamiento conversacional y GUI/herramientas, y mantiene la calidad mediante el baseline de SFT más los critical steps, de modo que las muestras de entrenamiento conservan en todo momento metadata interpretable.
